\documentclass[../../../../main.tex]{subfiles}
\begin{document}

自然数$m\geqq2$に対し，$m-1$個の二項係数
\begin{align*}
  _mC_1, \hspace{1\zw} _mC_2, \hspace{1\zw} \cdots, \hspace{1\zw} _mC_{m-1}
\end{align*}
を考え，これらすべての最大公約数を$d_m$とする．
すなわち$d_m$はこれらすべてを割り切る最大の自然数である．
\begin{enumerate}
  \item $m$が素数ならば，$d_m=m$であることを示せ．
  \item すべての自然数$k$に対し，$k^m-k$が$d_m$で割り切れることを，$k$に関する数学的帰納法によって示せ．
  \item $m$が偶数のとき$d_m$は1または2であることを示せ．
\end{enumerate}

\end{document}
